%% ============================================================
%% sections/a9-transparency.tex
%% H. R. 9510 Bill v3.0 - APPENDIX E. DEVELOPMENT TRANSPARENCY AND COMMIT SCHEDULE.
%% How Bill v3.0 was built. Rendered visuals: Table 9 (version lineage), Table 10
%% (single-PR commit schedule), Figure 11 (the build-and-handoff process).
%% ============================================================
\clearpage
\phantomsection\label{toc:a9}
\addcontentsline{toc}{section}{Appendix E. Development Transparency and Commit Schedule}
\usccrosshead{APPENDIX E. DEVELOPMENT TRANSPARENCY AND COMMIT SCHEDULE}

\uscprov{1}{\textbf{What was built, and how.} Bill v3.0 (documentation release
v3.1.0) is the finished visual amendment of H. R. 9510. It was produced by an
autonomous artificial intelligence coding agent, Claude Code Opus 4.8 (1M context,
Max), operating under human direction in a managed, ephemeral cloud container, by
executing the bracketed drafting instructions in the VVUQ-05 draft-bill scaffold
against the named figures-bill, next-steps, and VVUQ-04 final-bill sources. It was
authored across sequential commits in a single pull request, one file per commit
pushed to GitHub in real time, with a second-to-last error-fix-and-bundle commit
and a final repository-updates commit. Only kevinkawchak/cancer-automated was
edited \cite{anthropic-claude-code, repo-cancer-automated}.}

{\footnotesize
\begin{xltabular}{\textwidth}{@{}L{1.4cm} L{3.6cm} Y L{2.7cm}@{}}
\hline
\textbf{Version} & \textbf{Directory} & \textbf{Role} & \textbf{DOI}\\
\hline
\endfirsthead
\hline
\textbf{Version} & \textbf{Directory} & \textbf{Role} & \textbf{DOI}\\
\hline
\endhead
\hline
\endlastfoot
Bill v1.0 & papers/VVUQ-03/final-paper & Standalone Physical AI Oncology Trial
Bill & 10.5281/zenodo.20454870\\
\hline
Bill v2.0 & papers/VVUQ-04/final-bill & FD\&C Act amendment: new section 360e-5
and ten conforming amendments & 10.5281/zenodo.20485580\\
\hline
Bill v3.0 & papers/VVUQ-05/full-bill & This finished visual amendment &
10.5281/zenodo.xxxxxxxx\\
\hline
\end{xltabular}}
\figcaption{Table 9. Bill version lineage: each bill version, its directory, its
role, and its DOI.}

\uscprov{1}{\textbf{What this version adds.} From Bill v2.0, this version keeps the
full operative text and reorganizes it as a visual perspective: it renders eleven
text figures and ten full-width tables, adds the visual engineering evidence in
Appendix A, the submission deliverables in Appendix B, the explainability standard
in Appendix C, and the commit record here, while confining the emerging bills and
executive actions to Appendix D, exactly as the prior bill does
\cite{kawchak2026vvuq04bill}. The ordered commits of this single pull request are
listed in Table 10.}

{\footnotesize
\begin{xltabular}{\textwidth}{@{}L{1.6cm} L{4.6cm} Y@{}}
\hline
\textbf{Commits} & \textbf{Files or action} & \textbf{Role}\\
\hline
\endfirsthead
\hline
\textbf{Commits} & \textbf{Files or action} & \textbf{Role}\\
\hline
\endhead
\hline
\endlastfoot
1 to 3 & main.tex; usctitle.sty; references.bib & The assembler, the style, and
the sources\\
\hline
4 to 11 & sections/s2 through s4 and a5 through a9 & The bill body and the five
appendices, one file per commit\\
\hline
12 & README.md & The folder README (badges, tree, conventions)\\
\hline
13 to 25 & deliverables/01 through 12 and the package README & The submission
deliverables, processed as separate tasks, one commit each\\
\hline
26 & prompt-full-bill.md; output-full-bill.md; full-bill-LaTeX.zip & The
error-fix-and-bundle commit\\
\hline
27 & The main README, releases.md, CHANGELOG.md, and CITATION.cff & The final
repository-updates commit\\
\hline
\end{xltabular}}
\figcaption{Table 10. Single-pull-request commit schedule: the ordered commits
that build this bundle, one file per commit pushed in real time, ending with an
error-fix-and-bundle commit and a repository-updates commit.}

\uscprov{1}{\textbf{How this measure is implemented as law.} On enactment, the
Secretary issues final regulations within 365 days under section 515D(h) and may
issue interim guidance in the meantime; the amendments take effect on the first
day of the first calendar quarter beginning not less than one year after
enactment; an investigation already enrolling has 180 days to come into
compliance; a trial site must pass its Physical AI Standard Level gate and each
robot must meet its Unification Standard Level minimum before enrollment or
action; and, for each robot-patient interaction, the automated verification
process must clear, and the documentation and attestation must be filed, before
the code is generated or executed. The enforcement trigger is a single fact: the
absence of a cleared verification-before-generation record. Figure 11 shows the
build and the handoff to the deliverables package.}

\begin{asciifig}
BUILD AND HANDOFF (single pull request, real time)

  VVUQ-05/draft-bill (v3.0 scaffold: slots + bracketed instructions)
                          |
                          v
  THIS PASS  rendered every figure and table; wrote each deliverable
  +-----------------------------------------------------------------+
  | full-bill/  main.tex + .sty + .bib + 8 sections + README + zip  |
  | full-bill/deliverables/  12 documents + a package README        |
  +-----------------------------------------------------------------+
                          |
                          v
  the final visual bill (this PDF) and the submission-deliverables package
\end{asciifig}
\figcaption{Figure 11. The build-and-handoff process: the v3.0 scaffold rendered
in this pass into the finished visual bill and the deliverables package, all
within one pull request.}
